\chapter{La segunda salida}
% versión completa: chapters/16_chemoutput.tex

La máquina tiene dos salidas. La rápida: los músculos. La lenta: la química del cuerpo. El cerebro controla el corazón, los vasos sanguíneos y las glándulas, y a través de la sangre envía mensajes a todo el cuerpo a la vez. La sangre es el bus de difusión más amplio de todos: un mensaje llega a cada célula en más o menos un minuto, y su significado, como siempre, lo da la “cerradura” del receptor.

\section*{Acelerador y freno}

El corazón late solo; tiene su propio metrónomo. El cerebro solo ajusta el ritmo con dos cables de signo opuesto. \nt{NORA} es el pedal del acelerador: acelera, pero despacio. \nt{ACET} es el freno: frena rápido, porque no necesita una larga cadena de reacciones químicas. Y aquí por fin aparece trabajo para el tipo \nt{ADRE}: la adrenalina liberada a la sangre es el mismo pedal del acelerador, solo que pisado a la vez para todo el cuerpo.

\pencilfig{16}{\pencilart{16}}{Al corazón lo guían dos cables de signo opuesto: acelerador y freno.}

\section*{Una cascada con retroalimentación}

Muchas hormonas se controlan en cascada: el cerebro le da la orden a una primera glándula, esta a una segunda, la segunda libera la hormona, y la hormona le avisa al cerebro que ya hay suficiente. Es un termostato de tres etapas. Un ingeniero sabe qué peligro tiene un circuito así: si el lazo se hace demasiado “sensible”, empieza a oscilar, y el nivel, en vez de mantenerse, sube y baja. Por eso la precisión de un regulador así es limitada: no se puede hacer a la vez estable y muy preciso.

\pencilfig{127}{%
% three identical first-order stages with proportional feedback: secretion = max(0, K (target - level)).
% K = 3 settles at 3/4 of the target; K = 12 (above the stability limit K = 8) keeps oscillating. x = 0.42 t, y = 0.55 level
\foreach \yo/\t in {1.75/{lazo débil: estable, pero lejos de la meta}, 0/{lazo fuerte: más cerca de la meta, pero oscila}}{
  \draw[pen] (0,\yo) -- (8.5,\yo);
  \draw[pcGray, densely dashed, line width=0.5pt] (0,\yo+0.55) -- (8.5,\yo+0.55);
  \node[lbl, anchor=south west] at (2.0,\yo+0.8) {\t};}
\node[lbl, anchor=west] at (8.55,2.3) {meta}; \node[lbl, anchor=west] at (8.55,0.55) {meta};
\begin{scope}[yshift=1.75cm]
\draw[penlite=pcBlue] plot[smooth] coordinates {(0.00,0.00) (0.17,0.01) (0.34,0.08) (0.50,0.19) (0.67,0.33) (0.84,0.46) (1.01,0.55) (1.18,0.59) (1.34,0.58) (1.51,0.55) (1.68,0.49) (1.85,0.43) (2.02,0.37) (2.18,0.34) (2.35,0.33) (2.52,0.34) (2.69,0.37) (2.86,0.40) (3.02,0.43) (3.19,0.44) (3.36,0.45) (3.53,0.45) (3.70,0.44) (3.86,0.42) (4.03,0.41) (4.20,0.40) (4.37,0.39) (4.54,0.39) (4.70,0.40) (4.87,0.40) (5.04,0.41) (5.38,0.42) (5.88,0.42) (6.38,0.41) (7.06,0.41) (8.40,0.41)};
\end{scope}
\draw[penlite=pcRed] plot[smooth] coordinates {(0.00,0.00) (0.17,0.05) (0.34,0.30) (0.50,0.68) (0.67,1.02) (0.84,1.23) (1.01,1.30) (1.18,1.26) (1.34,1.16) (1.51,1.02) (1.68,0.87) (1.85,0.72) (2.02,0.58) (2.18,0.47) (2.35,0.39) (2.52,0.38) (2.69,0.46) (2.86,0.57) (3.02,0.67) (3.19,0.71) (3.36,0.70) (3.53,0.65) (3.70,0.58) (3.86,0.50) (4.03,0.43) (4.20,0.41) (4.37,0.45) (4.54,0.53) (4.70,0.61) (4.87,0.65) (5.04,0.64) (5.21,0.60) (5.38,0.54) (5.54,0.47) (5.71,0.42) (5.88,0.42) (6.05,0.48) (6.22,0.56) (6.38,0.62) (6.55,0.64) (6.72,0.62) (6.89,0.56) (7.06,0.50) (7.22,0.44) (7.39,0.42) (7.56,0.45) (7.73,0.53) (7.90,0.60) (8.06,0.63) (8.23,0.62) (8.40,0.58)};
\node[lbl, anchor=north] at (4.2,-0.05) {tiempo};
}{Modelo de una cascada hormonal de tres etapas. Una retroalimentación débil calma el nivel, pero lo deja bastante por debajo de la meta. Una fuerte lo acerca más, pero entonces el nivel ya no se mantiene, sino que oscila.}

\section*{Pulsos, no nivel}

Algunas hormonas solo funcionan en pulsos. Si se administran de forma continua, la glándula receptora “se cansa” y se apaga; si llegan en pulsos breves una vez por hora, funciona. El receptor no lee el nivel, sino el ritmo, como la sinapsis del capítulo del código Morse que solo oye los cambios.

\section*{El día}

Y por último, en el cuerpo hay un generador lento con un período de alrededor de un día. Su período natural difiere un poco de veinticuatro horas, y cada mañana la luz lo ajusta, como un receptor de radio se sintoniza con una emisora. No lo confunda con el generador de reloj de una computadora: no mide los pasos del cálculo. Cambia los modos de trabajo de toda la máquina: día y noche, vigilia y sueño.

\fullversion{16}
